\documentclass{article}

\usepackage{amsmath,amssymb}
\usepackage{xurl}
\usepackage[colorlinks=true,linkcolor=blue,citecolor=blue,urlcolor=blue]{hyperref}
\setlength{\emergencystretch}{2em}

% Shared commands for notes in docs/paper-gaps/.

\DeclareUnicodeCharacter{2080}{\ifmmode{}_{0}\else\textsubscript{0}\fi}
\DeclareUnicodeCharacter{2081}{\ifmmode{}_{1}\else\textsubscript{1}\fi}
\DeclareUnicodeCharacter{2082}{\ifmmode{}_{2}\else\textsubscript{2}\fi}
\DeclareUnicodeCharacter{2099}{\ifmmode{}_{n}\else\textsubscript{n}\fi}

\newcommand{\Lean}{\textsc{Lean}}
\ExplSyntaxOn
\NewDocumentCommand{\leanid}{m}{
  \ifmmode
    \textsf{\detokenize{#1}}
  \else
    \group_begin:
    \urlstyle{sf}
    \tl_set:Nn \l_tmpa_tl {#1}
    \tl_replace_all:Nnn \l_tmpa_tl {\_} {_}
    \exp_args:NV \path \l_tmpa_tl
    \group_end:
  \fi
}
\ExplSyntaxOff
\providecommand{\tr}{\operatorname{tr}}
\newcommand{\tnleanIssueBase}{https://github.com/LionSR/TNLean/issues/}
\newcommand{\tnleanPrBase}{https://github.com/LionSR/TNLean/pull/}
\newcommand{\tnleanDocBase}{https://sirui-lu.com/TNLean/docs/}
\newcommand{\ghissue}[1]{\href{\tnleanIssueBase#1}{GitHub issue \##1}}
\newcommand{\ghpr}[1]{\href{\tnleanPrBase#1}{PR \##1}}
\newcommand{\doclink}[1]{\href{\tnleanDocBase#1.html}{\path{#1}}}

% Dirac notation and the identity operator, as in blueprint/src/macros/common.tex.
% \providecommand keeps note-local definitions authoritative.
\usepackage{dsfont}
\providecommand{\ket}[1]{|#1\rangle}
\providecommand{\bra}[1]{\langle #1|}
\providecommand{\braket}[2]{\langle #1 | #2 \rangle}
\providecommand{\ketbra}[2]{|#1\rangle\!\langle #2|}
\providecommand{\Id}{\mathds{1}}
\providecommand{\one}{\mathds{1}}
\providecommand{\C}{\mathbb{C}}
\providecommand{\MN}[1]{M_{#1}(\C)}
\providecommand{\MD}{M_D(\C)}

% Verdict marker: \gapnote{<kind>}{<status>}. Kind is the policy
% classification (clarification, local-correction, scope-restriction,
% unfaithful, false-source, open-gap); status is open, wip (elimination
% underway), resolved, or historical. Machine-read by the site index;
% prints a verdict line.
\newcommand{\gapnote}[2]{\par\noindent\textbf{Verdict:} #1 (#2).\par}


% Fallbacks keep this author document compilable when it is built outside its
% source directory and a different command.tex is found first.
\providecommand{\leanid}[1]{\textsf{\detokenize{#1}}}
\providecommand{\doclink}[1]{%
    \href{https://sirui-lu.com/TNLean/docs/#1.html}{\path{#1}}}
\providecommand{\gapnote}[2]{%
    \par\noindent\textbf{Verdict:} #1 (#2).\par}

\title{The Consequences of the Distribution Protocol Used in the Assembly}
\author{}
\date{2026-10-17}

\begin{document}
\maketitle
\gapnote{scope-restriction}{open}

\section{The source proposition}

Proposition~7.1 of the polynomial PEPS approximation manuscript of
September~24, 2026~\cite{OpenAI2026PolynomialPEPS} (\path{06-geometry.tex},
lines~12--53, at revision \texttt{adc7f1241b42e322a6451854ab7e4b4c146bf78a})
asserts that for all sufficiently large grid side lengths $L$, with constants
uniform in the Hamiltonian and in $L$, there exists a protocol that distributes
the ground vector $\Omega$ through the padded dyadic hierarchy. Its four items
describe:
\begin{enumerate}
\item the parties and their product initialization;
\item the division of the chronology into changes, their gates, and the bounded
  lifetime participation of each party;
\item the dyadic scales and grid addresses of the parties acted on by each
  gate;
\item the sequence of intermediate reference vectors, ending at a tensor
  product $\Omega \otimes \Xi$ with an auxiliary vector $\Xi$.
\end{enumerate}

\section{The formal statement}

The formal statement defines the subset of properties of Proposition~7.1 that
the PEPS assembly of Section~8 uses. This retained statement is proved from
the adaptive patch constraint of Proposition~4.1 of
Ref.~\cite{OpenAI2026PolynomialPEPS}, which is assumed; the angular information
bound of Proposition~3.4 is proved.\footnote{The theorem is
    \leanid{TNLean.PEPS.Approximation.largeSizeDyadicProtocolConsequences_of_adaptivePatchConstraint}.}
The full Proposition~7.1 remains open until the omitted clauses listed below
are proved as well. A protocol is formalized as a distributed
construction in the sense of Theorem~5.2 of Ref.~\cite{OpenAI2026PolynomialPEPS}
whose chronology is divided into $n$ changes, with each party attached to a
square of the padded dyadic hierarchy and with a reference vector $r_t$ at each
stage $0 \le t \le n$, starting at the initial product vector $r_0$. Let
$C > 0$ be a scale coefficient and let $b, k, r \in \mathbb N$ be fixed degree
and exponent bounds. The retained properties are:
\begin{enumerate}
\item there are at most $C L^2$ parties in total, and at most $b$ parties are
  attached to any square of the padded hierarchy; a different party is
  assigned to each aligned square of the padded hierarchy, including the root
  square, at the scale and address of that square (\path{06-geometry.tex},
  lines~18--20 and~64--67); every other party is an additional party with the
  scale and address of a non-root square of the hierarchy, and it takes part
  only in changes of the repainting of that square (lines~19--22 and~35--37).
  Each change records the square of the repainting it belongs to, if any; the
  changes of the repainting of a square form one contiguous segment of the
  schedule, and every party acted on in such a change has the scale of the
  square or an adjacent one and an address within $b$ times the side of the
  square of its address. The source bounds the additional parties of each
  repainting; the retained bound on the parties at each square implies this,
  since the additional parties of a repainting are placed at its square;
\item initialization is a pure product of copies of $\Omega$: the initial
  encoded frames are the main frame followed by the auxiliary frames, each
  without holes and held entirely by one party, and the initial memory and
  vector of each party are those of the copies it holds (\path{06-geometry.tex},
  lines~21--23); the initial product vector $r_0$ is carried by private maps
  alone to the reference vector of these frames;
\item there are at most $C L^3$ changes, each change contains at most $b$
  nonprivate gates acting on at most $b$ parties each, and each party
  participates in at most $b$ changes over its lifetime;
\item each gate is a sum of at most $C L^k$ monomials with absolute coefficient
  sum at most $C L^k$ and at most $r$ pair effects per monomial;
\item any two parties acted on by a gate have equal or adjacent dyadic scales
  and grid addresses separated by at most $b$ times either endpoint scale;
\item if $W_t$ denotes the contraction of stage $t$, the reference vectors
  $r_t$ and terminal auxiliary vector $\Xi$ satisfy
  \begin{align}
      \lVert W_t r_t - r_{t+1}\rVert
          &\le C L^{-20},
      \label{eq:prop71_step_error}\\
      r_n
          &= \Omega \otimes \Xi,
      \label{eq:prop71_terminal_factorization}\\
      1 - C L^{-57}
          &\le \lVert \Xi\rVert \le 1.
      \label{eq:prop71_junk_norm}
  \end{align}
\item every intermediate reference vector $r_t$ is the product of the
  reference vectors of the encoded frames of $\Omega$ at that stage, the main
  encoded frame and the auxiliary encoded frames, all of norm at most one
  (\path{06-geometry.tex}, lines~38--40); before the readout the raw registers
  of the main frame are held by the parties of the genuine sites. The main
  frame is singled out at the initial stage and before the readout; at the
  stages in between, the order of the factors is not recorded. The last stage
  is the readout of the physical and discarded registers, which has no
  counterpart in the source and is described by
  (\ref{eq:prop71_terminal_factorization}).
\end{enumerate}
From these properties and (\ref{eq:prop71_step_error})--(\ref{eq:prop71_junk_norm}),
the formal development proves the gate-density estimate of Section~8.1, the
link-separation hypothesis of Lemma~8.2, and the lifetime bound of at most
$64 G$ gates per block party when each change has at most $G$
gates.\footnote{The protocol predicate is
    \leanid{TNLean.PEPS.Approximation.DyadicProtocol.IsDistributionProtocol};
    the large-size existence hypothesis is
    \leanid{TNLean.PEPS.Approximation.LargeSizeDyadicProtocolConsequences}.}

\section{Omitted clauses}

Two clauses of Proposition~7.1 are omitted from the formal predicate:
\begin{enumerate}
\item Item~3 of Theorem~5.2 (\path{04-compression.tex}, lines~143--147) also
  bounds the pair sources and the messages of each monomial before expansion.
  The formal statement counts only pair effects.
\item The paper provides expansions accurate to any prescribed inverse power of
  $L$, whereas (\ref{eq:prop71_step_error}) fixes a single contraction $W_t$
  per change with error $C L^{-20}$.
\end{enumerate}
Neither of these two clauses is needed for the Section~8 assembly: the gate-density
bound, the routing separation hypothesis, and the lifetime gate count depend
only on the retained bounds, while the message and replacement-slot bounds used
in compression follow from the message expansion and the nonprivate gate count
and arity. The formal statement is therefore a weaker consequence of
Proposition~7.1 sufficient for Section~8, proved modulo Proposition~4.1, while
the full Proposition~7.1 remains open.

\section{Elimination}

Formalize the full statement of Proposition~7.1 including the pre-expansion
counts of pair sources and messages for arbitrary inverse-power accuracy.
Proving Proposition~7.1 in that stronger form then implies the present
existence hypothesis by projection onto the retained fields.

\bibliographystyle{alpha}
\bibliography{references}

\end{document}
